\section{行动清单与下一步}

\subsection{部署前准备}

在 rollout 之前，需要收敛以下交付物：
\begin{itemize}
    \item \textbf{Data readiness}：确认 preference dataset 覆盖新的 prompt cluster，避免 distribution shift。
    \item \textbf{Model readiness}：检查 RLHF 有 limit order，DPO/dRL 版本保持 KL margin。
    \item \textbf{Governance readiness}：所有 guardrail policies 都写入 policy doc，compliance team 签署。
\end{itemize}

\begin{knowledgebox}{准备阶段的自动化清单}
把 readiness checklist 反向链接到 issue tracker，每个 item 通过 automation bot ping 相关 owner，确保没有漏项。
\end{knowledgebox}

\subsection{上线后反馈}

部署后要保持持续反馈循环：
\begin{itemize}
    \item 依赖 \textbf{evidence dashboard} 触发 auto-alert，当 reward drift+latency spikes 同时发生时自动 escalate。
    \item 每周汇报 \textbf{ops scorecard}：reward score drift、human override、rejection trend。
    \item 用 \textbf{post-deployment study} 检查 new prompts、adversarial tests、user complaints 是否同步反馈到 labeling pipeline。
\end{itemize}

\begin{importantbox}{持续反馈的重要性}
Auto-alert 帮助快速发现 drift，但真正的 alignment 还是要靠 weekly review meeting，把 quantitative signals 转换成 qualitative insights。
\end{importantbox}

\subsection{本章小结}

完整的行动清单包括部署前的 readiness checklist 与 deployment 后的 feedback loop，只有严密闭环才能把 RLHF pipeline 的 research 变成 resilient product。

